利用同构 \eqref{E:gr as quotient}，容易看出 $\Gr\tilde\times\Gr\cong \Gr^{(\infty)}\times^K\Gr=LG\times^K\Gr$，
其卷积映射由作用映射 $LG\times\Gr\to \Gr$ 诱导。
这促使我们定义 $\Gr$ 的 $K$-不变闭子集的扭积。即：
对闭子集 $X\subset \Gr$，令 $X^{(\infty)}$ 表示它在 $\Gr^{(\infty)}$ 中的原像。则 $X^{(\infty)}\to X$ 是 $K$-torsor；对每个 $n$，令 $X^{(n)}=X^{(\infty)}/K^{(n)}\subset LG/K^{(n)}$，它是 $X$ 上的 $K_n$-torsor。
于是给定两个闭子集 $X_1, X_2$，其中 $X_2$ 是 $K$-不变的，我们把它们的扭积记作
\begin{equation}
\label{E-conv Sch}
X_1\tilde\times X_2:= X_1^{(\infty)}\times^K X_2,
\end{equation}
它是 $\Gr\tilde\times\Gr$ 的闭子集（因而可表）。若 $X_2$ 拟紧，则存在充分大的整数 $m$，使 $K$ 在 $X_2$ 上的作用经 $K_m$ 分解。
此时 $X_1^{(\infty)}\times^K X_2=X_1^{(m)}\times^{K_m}X_2$。注意
若 $X_i=\Gr_{\mu_i}$（$\mu_i\in \xcoch(T)/W$），我们也可把 $\Gr_{\mu_1}\tilde{\times}\Gr_{\mu_2}$ 描述为
\[\Gr_{\mu_1}\tilde{\times}\Gr_{\mu_2}\cong \left\{(\mE_1,\mE_2,\beta_1,\beta_2)\in\Gr\tilde\times\Gr\mid \inv(\beta_1)\leq \mu_1, \inv(\beta_2)\leq \mu_2\right\}.\]

更一般地，对 $r$ 重卷积 Grassmann 簇中的闭子集 $X_1$ 与 $s$ 重卷积 Grassmann 簇中的 $K$-不变闭子集 $X_2$，可用 \eqref{E-conv Sch} 中同一公式定义它们的扭积，其中 $X_1^{(\infty)}$ 理解为 $X_1$ 在 $(\Gr\tilde\times\cdots\tilde\times\Gr)^{(\infty)}$ 中的原像。
特别地，若 $\mmu=(\mu_1,\ldots,\mu_r)$ 是 $G$ 的支配余权序列，可定义 Schubert 簇的扭积
$$\Gr_{\mmu}:=\Gr_{\mu_1}\tilde\times\cdots\tilde\times\Gr_{\mu_r}\subset\Gr\tilde\times\cdots\tilde\times\Gr,$$
而若 $\nu_\bullet=(\nu_1,\ldots,\nu_s)$ 是另一序列，则 $\Gr_{\mmu}\tilde\times\Gr_{\nu_\bullet}\cong\Gr_{\mmu,\nu_\bullet}$。
令 $|\mu_\bullet|=\sum \mu_i$。则卷积映射 \eqref{conv m} 诱导
\begin{equation}\label{conv m2}
m: \Gr_{\mu_\bullet}\to \Gr_{|\mu_\bullet|}, \quad \mE_r\stackrel{\beta_r}{\dashrightarrow}\cdots\stackrel{\beta_0}{\dashrightarrow}\mE_0=\mE^0\mapsto (\mE_r,\beta_1\cdots\beta_r),
\end{equation}
已知它是半小映射。这意味着
\begin{equation}
\label{E:dim est}
\dim (\Gr_{\mmu}\times_\Gr\Gr_{\nu_\bullet})\leq \frac{1}{2}(d_{\mmu}+d_{\nu_\bullet}),
\end{equation}
其中记 $d_{\mmu}=\dim \Gr_{\mmu}$（见 \cite[命题 3.1.10]{XZ}）。

现在，对 $\mA_1,\mA_2\in \on{P}_{K}(\Gr)$，
以 $\mA_1\tilde\boxtimes\mA_2$ 记 $\mA_1$ 与 $\mA_2$ 在 $\Gr\tilde\times\Gr$ 上的「外扭积」。它是 $\Gr\tilde\times \Gr$ 上唯一的反常层，使其到 $\on{Supp}(\mA_1)^{(m)}\times \on{Supp}(\mA_2)$ 的拉回同构于 $\mA_1$ 到 $\on{Supp}(\mA_1)^{(m)}$ 的拉回与 $\mA_2$ 的外积。这里 $\on{Supp}(\mA_i)$ 是 $\mA_i$ 的支撑，它是 $K$-不变的；$m$ 充分大，使 $K$ 在 $\on{Supp}(\mA_2)$ 上的作用经 $K_m$ 分解。
例如若 $\mA_i=\IC_{\mu_i}$，则 $\IC_{\mu_1}\tilde\boxtimes\IC_{\mu_2}$ 典范同构于 $\Gr_{\mu_1}\tilde\times\Gr_{\mu_2}$ 的相交上同调层。
于是 $\mA_1$ 与 $\mA_2$ 的卷积积定义为
\begin{equation}\label{Lus conv}
\mA_1\star\mA_2:=m_!(\mA_1\tilde\boxtimes\mA_2),
\end{equation}
其中 $m:\Gr\tilde\times\Gr\to \Gr$ 是卷积映射（由 \eqref{conv m0} 定义）。
先验地，这是 $\Gr$ 上的 $K$-等变复形；但由卷积映射 \eqref{conv m2} 的半小性，它是反常层。

扭积 $\Gr\tilde\times\Gr$ 「拓扑同构于」乘积 $\Gr\times\Gr$。若 $F=\bC((\varpi))$，可用 $\Gr$ 上的解析拓扑使这句话精确化。即在解析拓扑下，$K$-torsor $\Gr^{(\infty)}\to \Gr$ 平凡。一般情形可从动机上理解。
于是当 $E=\Ql$（$\ell\neq \cha k$）时，应用 K\"unneth 公式，可给超上同调函子
$$\on{H}^*(\Gr_G,-) : \on{P}_{K}(\Gr_G,\Ql)\to \on{Vect}_{\Ql}$$ 配上典范幺半结构（见 \cite[\S 2.3]{Zh1} 与 \cite[\S 5.2]{Zh2}）。

下面这个定理通常称为几何 Satake 等价。我们假定 $E=\Ql$。

\begin{thm}\label{intro:geomSat}
幺半函子 $\on{H}^*$ 分解为两个函子的复合：一个是从 $\Sat_G$ 到 Langlands 对偶群 $\hat{G}$ 在 $\Ql$ 上有限维表示范畴 $\on{Rep}_{\Ql}(\hat{G})$ 的幺半范畴等价，另一个是从 $\on{Rep}_{\Ql}(\hat{G})$ 到 $\Ql$ 上有限维向量空间范畴 $\on{Vect}_{\Ql}$ 的遗忘函子。在该等价下，Schubert 簇 $\Gr_\mu$ 的相交上同调层 $\IC_\mu$ 对应于 $\hat G$ 的「最高权」为 $\mu$ 的不可约表示 $V_\mu$。
\end{thm}


\begin{rmk}
\label{R: birth of the dual group}
(1) 事实上，更典范的做法是把 $G$ 的对偶群 $\hat G$ \textbf{定义}为 Tannakian 范畴 $(\Sat_G,\on{H}^*)$ 的 Tannakian 群，并把 $V_\mu$ \textbf{定义}为 $\on{H}^*(\Gr,\IC_\mu)$，配以重言作用 $\hat G=\Aut^\otimes(\on{H}^*)$。本文余下部分采取这一观点。

(2) 如 \cite[\S 5.3]{Zh2} 所述，Tannakian 群典范地配有带钉 $(\hat G, \hat B, \hat T, \hat X)$。我们简要回忆 $(\hat B, \hat T)$ 的构造，细节参看「同上」。
上同调 $\on{H}^*$ 上的分次给出余特征标 $\bG_m\to \hat G$，可以证明它是正则的。于是它的中心化子给出极大环面 $\hat T$，而 $\hat B\supset \hat T$ 是 $\hat G$ 中使该余特征标为支配的唯一点 Borel 子群。
\end{rmk}

当 $F=k((t))$ 是等特征局部域时，此定理是 Lusztig、Ginzburg、Beilinson–Drinfeld 与 Mirkovi\'{c}-Vilonen 等人工作的结果（参看 \cite{Lu,Gi1,BD,MV}）。当 $F$ 是混合特征局部域时，此定理在 \cite{Zh1} 中证明。

定理最难的部分，是在 $\on{P}_{K}(\Gr_G)$ 的幺半结构上构造交换性约束，使 $\on{H}^*$ 成为张量函子。在等特征情形，该构造依赖把卷积积解释为融合积，这关键地用到 $k$ 上代数曲线笛卡尔幂上的 Beilinson–Drinfeld Grassmann 簇的存在性。如注 \ref{R:scholzediamond} 所述，混合特征情形中类似的几何对象已由 Scholze 近期构造。Scholze 近期还宣布，可以用与等特征情形类似的方式，在混合特征情形定义融合积并建立几何 Satake。

\cite{Zh1} 所用的方法不同。我们用经典 Gelfand 技巧的范畴版构造交换性约束。这一想法并不新，已见于 \cite{Gi1}。因此即使在混合特征情形，我们确实有交换性约束的候选。问题在于：若不使用融合解释，就难以验证它应当满足的性质（例如六边形公理）。

我们的新观察是：这些性质成立，等价于仿射 Hecke 代数的一个数值结果。即在 \cite{LV,Lu2} 中，Lusztig 与 Vogan 对带对合的 Coxeter 系 $(W,S)$ 引入了某些多项式 $P^\sigma_{y,w}(q)$，它们类似于通常的 Kazhdan–Lusztig 多项式 $P_{y,w}(q)$。随后 \cite{Lu2} 中猜想：若 $(W,S)$ 是仿射 Weyl 群，$y,w$ 是 $W$ 中某些元素，
\[P_{y,w}^\sigma(q)=P_{y,w}(-q).\]
